\section*{Chapitre \ref{ch:dv:rough-volatility-and-forward-variance-models} --- Volatilité rugueuse et modèles de variance forward}

\begin{solution}{exo:dv:rough-volatility-and-forward-variance-models:1}
Les variances totales valent $0.04$ et $2\times0.0484=0.0968$~; la \omterm{def:dv:rough-volatility-and-forward-variance-models:fwdvar}{variance forward} vaut $0.0568$ par an, soit une volatilité forward
de 23.8~\%.
\end{solution}

\begin{solution}{exo:dv:rough-volatility-and-forward-variance-models:2}
Le strip pondère chaque option hors de la monnaie par $1/K^2$, de sorte que les bas strikes comptent davantage que les hauts. Avec un
skew baissier, les puts à bas strike sont cotés à des volatilités plus élevées que l'option à la monnaie, et la
moyenne pondérée des variances est supérieure à la variance à la monnaie.
\end{solution}

\begin{solution}{exo:dv:rough-volatility-and-forward-variance-models:3}
$m(2,\Delta)\propto\Delta^{2H}$, donc $2^{2H}=1.15$ et $H=\ln1.15/(2\ln2)=0.10$.
\end{solution}

\begin{solution}{exo:dv:rough-volatility-and-forward-variance-models:4}
$Y_t$ est gaussien, de moyenne nulle et de variance $2H\int_0^t(t-s)^{2H-1}ds=t^{2H}$, donc
$\E[e^{\eta Y_t}]=e^{\eta^2t^{2H}/2}$ et $\E[v_t]=\xi_0(t)$. Pour $u\ge t$,
$\xi_t(u)=\E_t[v_u]$~: on décompose $Y_u$ en $\sqrt{2H}\int_0^t(u-s)^{H-1/2}dW_s$, connu en $t$, et un
reste indépendant de variance $(u-t)^{2H}$~; en prenant l'espérance de son exponentielle, on retrouve la formule du
texte. Sa seule dépendance en $t$ passant par $W$ est $\eta\sqrt{2H}\int_0^t(u-s)^{H-1/2}dW_s$, dont la différentielle est
$\eta\sqrt{2H}(u-t)^{H-1/2}dW_t$~: c'est la volatilité de $\xi_t(u)$.
\end{solution}

\begin{solution}{exo:dv:rough-volatility-and-forward-variance-models:5}
Le contrat à terme vaut $\E[\mathrm{VIX}]$ et la volatilité forward $\sqrt{\E[\mathrm{VIX}^2]}$~; par l'inégalité de Jensen
$\E[\mathrm{VIX}]\le\sqrt{\E[\mathrm{VIX}^2]}$, avec un écart qui croît avec la variance du VIX. Une valeur
plus grande de $\eta$ rend le VIX plus volatil et élargit l'écart~: avec $\eta=2.5$, le contrat à terme vaut 17.88.
\end{solution}

\begin{solution}{exo:dv:rough-volatility-and-forward-variance-models:6}
Dans l'état calme $R_1=0$ et $R_2=\sigma^2$, donc $\sigma^*=\beta_0+\beta_2\sigma^*$ et $\sigma^*=0.04/0.35=11.4\%$.
Après une journée à $+4\%$, $\beta_1R_1$ est négatif et plus grand en valeur absolue que la hausse de $\beta_2\sqrt{R_2}$, donc
la volatilité baisse~; $R_1$ décroît au taux 25 et $R_2$ au taux 15, donc le terme d'activité survit au terme de tendance et
la volatilité remonte au-dessus de $\sigma^*$, jusqu'à 12.4~\%, avant que les deux termes ne s'estompent.
\end{solution}

\begin{solution}{exo:dv:rough-volatility-and-forward-variance-models:7}
Sur cinq graines~: de 0.45 à 0.50 sans bruit, le $\frac12$ de la diffusion, légèrement abaissé par le retour à la moyenne~; 0.20
à 0.23 avec bruit. Le chemin rugueux du même code affiche 0.10 à 0.12.
\end{solution}

\begin{solution}{exo:dv:rough-volatility-and-forward-variance-models:8}
La volatilité réalisée journalière est mesurée avec une erreur d'environ 0.08 en log, indépendante d'un jour à l'autre. L'erreur
ajoute une constante à la fonction de structure à chaque décalage et en aplatit la pente~: une diffusion de $H=\frac12$
lue à travers elle donne environ 0.17 à 0.23. L'estimation ne réfute rien tant que le bruit n'est pas modélisé, par exemple en retranchant sa contribution, deux fois
sa variance, de $m(2,\Delta)$ avant la régression~; l'indication
fournie par la structure par terme du skew, où aucun bruit de mesure n'intervient, est le test le plus solide.
\end{solution}

\begin{solution}{pb:dv:rough-volatility-and-forward-variance-models:1}
\textbf{1.} 22.1~\% et 19.2~\%.
\textbf{2.} 24.0~\%.
\textbf{3.} 25.8~\%.
\textbf{4.} 24.9~\%.
\textbf{5.} Oui~: la variance totale croît à chaque pilier, et \texttt{negative\_intervals()} est vide.
\textbf{6.} 4.13, 2.30 et 0.85.
\textbf{7.} $\eta=2.51$, $\kappa=1.08$.
\textbf{8.} 3.1 (7.76 contre 2.50).
\textbf{9.} L'exponentielle vaut au plus $\eta$ à horizon nul~; la loi de puissance y est infinie. Deux points fixent
$\eta$ et $\kappa$, et les courbes diffèrent ensuite partout ailleurs.
\textbf{10.} Elle ajoute un terme à décroissance rapide qui relève le court terme~; une somme d'exponentielles peut imiter la loi de
puissance sur une plage d'horizons choisie, jamais jusqu'à zéro.
\textbf{11.} $-1.44$ et $-0.25$.
\textbf{12.} $-0.70$ et $-0.25$.
\textbf{13.} $-0.44$ pour les deux.
\textbf{14.} $-0.09$~: presque plate.
\textbf{15.} $1.44/0.25=5.7$, et une loi de puissance d'exposant $-0.44$ prédit $52^{0.44}=5.7$ entre une semaine et un
an.
\textbf{16.} $H=\frac12-0.44=0.06$.
\textbf{17.} La règle appliquée à un modèle de $H=0.1$ connu renvoie aussi 0.06~: sur cette plage d'échéances, elle
sous-estime d'environ 0.04, donc la surface est compatible avec $H$ voisin de 0.1.
\textbf{18.} Les séries temporelles donnent 0.075 à 0.158 sur 21 indices, 0.124 pour le S\&P~500 dans le même tableau~: compatible
avec 0.1.
\textbf{19.} Exposant $-0.44$~; $H=0.06$ selon la règle des échéances courtes, environ 0.1 une fois la règle calibrée sur le
\omterm{def:dv:rough-volatility-and-forward-variance-models:rbergomi}{modèle de Bergomi rugueux}.
\textbf{20.} Oui pour les options sur indice à courte échéance, dont Heston ne peut pas produire le skew. Coûts~: Monte-Carlo seulement (pas d'EDP,
\omterm{def:dv:greeks-and-the-hedging-pnl:greeks}{grecques} plus lentes et plus bruitées), et le smile du VIX que le modèle implique doit être confronté à celui du marché avant que ses
produits sur la \omterm{def:dv:stochastic-volatility:sv}{volatilité de la volatilité} ne soient jugés fiables.
\end{solution}

\begin{solution}{iq:dv:rough-volatility-and-forward-variance-models:1}
$\xi_t(u)=\E^{\mathbb Q}_t[v_u]$, la variance instantanée espérée à une date future. La variance intégrée espérée
jusqu'à $T$ vaut moins deux fois le prix du contrat log, que répliquent un strip d'options hors de la monnaie~; c'est le
strike du swap de variance $T\sigma^2_{\mathrm{VS}}(T)$, et sa dérivée en $T$ vaut $\xi_0(T)$.
\iqlookfor{la définition, le strip du contrat log, et la dérivée par rapport à l'échéance.}
\end{solution}

\begin{solution}{iq:dv:rough-volatility-and-forward-variance-models:2}
Le logarithme de la volatilité se comporte aux courtes échelles de temps comme un mouvement brownien fractionnaire d'exposant de Hurst $H$ voisin de 0.1, bien
plus rugueux que le $\frac12$ d'une diffusion. Preuves~: la loi d'échelle de la fonction de structure de la volatilité réalisée
sur de nombreux indices et moments, et le skew à la monnaie en loi de puissance $\psi\sim T^{H-1/2}$ dans les prix d'options. Mise en garde~: le
temps, toute estimation sur séries temporelles est sensible au bruit de mesure.
\iqlookfor{les deux types de preuves et la mise en garde.}
\end{solution}

\begin{solution}{iq:dv:rough-volatility-and-forward-variance-models:3}
Dans un modèle de volatilité stochastique diffusif, le skew court terme provient de la volatilité de la variance sur un court laps de temps,
qui est borné, donc le skew tend vers une constante lorsque $T\to0$. Une \omterm{def:dv:rough-volatility-and-forward-variance-models:rough}{volatilité rugueuse} a une volatilité de la
\omterm{def:dv:rough-volatility-and-forward-variance-models:fwdvar}{variance forward} qui explose comme $\tau^{H-1/2}$ aux courts horizons, produisant un skew $\sim T^{H-1/2}$ qui continue
de croître, comme celui du marché.
\iqlookfor{la limite aux échéances courtes, et le noyau.}
\end{solution}

\begin{solution}{iq:dv:rough-volatility-and-forward-variance-models:4}
Le modèle n'est pas markovien, donc Monte-Carlo. Exact~: le processus de Volterra et son mouvement brownien sont conjointement gaussiens sur
la grille, donc factorisation de Cholesky de la covariance $2n\times2n$, en $O(n^3)$ une fois et $O(n^2)$ par trajectoire. Plus rapide~: le schéma hybride,
exact près de la singularité et une convolution par transformée de Fourier rapide ailleurs, $O(n\log n)$. Valoriser avec la \omterm{prop:dv:stochastic-volatility:mixing}{formule de mélange}
(en conditionnant sur le moteur de variance) pour réduire le bruit~; il faut vérifier la propriété de martingale.
\iqlookfor{non-markovien, la structure gaussienne, les coûts, la réduction de variance.}
\end{solution}

\begin{solution}{iq:dv:rough-volatility-and-forward-variance-models:5}
Calibrer à la fois le smile de l'indice et le smile du VIX avec un seul modèle. Le smile de l'indice contraint le skew et la volatilité de
la variance à courte échéance~; le smile du VIX contraint directement la loi de cette variance, et les modèles simples ne peuvent pas faire les deux
(conjecture de Guyon~; à laquelle répondent le \omterm{def:dv:stochastic-volatility:heston}{modèle de Heston} rugueux quadratique, le \omterm{def:dv:rough-volatility-and-forward-variance-models:pdv}{modèle de volatilité dépendante du chemin} et la construction exacte
de Guyon). Un desk s'y intéresse, car les produits sur le VIX et les produits sur indice sensibles à la \omterm{def:dv:stochastic-volatility:sv}{volatilité de la volatilité}
(cliquets, options sur variance) sont valorisés sur exactement la dynamique qu'un calage unilatéral laisse non testée.
\iqlookfor{ce que contraint chaque smile, et les produits exposés au risque.}
\end{solution}

\begin{solution}{iq:dv:rough-volatility-and-forward-variance-models:6}
Ils s'accordent sur les vanilles aux échéances cotées et divergent sur tout ce qui dépend de la volatilité de la
\omterm{def:dv:rough-volatility-and-forward-variance-models:fwdvar}{variance forward} à d'autres horizons~: options à courte échéance en deçà des échéances cotées, options à départ différé et cliquets
(skew forward), options sur variance et sur le VIX. Il faut trancher en testant chacun face aux instruments qui mesurent
ces dynamiques (options hebdomadaires, contrats à terme et options sur le VIX, cotations à départ différé) et en comparant le
\omterm{def:dv:greeks-and-the-hedging-pnl:hpnl}{P\&L de couverture} que chaque modèle aurait produit~; il faut conserver une réserve de modèle pour l'écart sur les produits qui restent.
\iqlookfor{la liste des produits, un test face à d'autres marchés, et une réserve de modèle.}
\end{solution}
